\documentclass[12pt]{article}
\usepackage[a3paper, margin=1in]{geometry}
\usepackage{amsmath, amssymb, amsthm, ulem, booktabs}

\begin{document}

\section*{Domácí úkol 11.1}

Vypracoval: Daniel \textit{``Randál"} Ransdorf \hfill Podpis: \rule{4cm}{0.4pt}

\begin{enumerate}
  
  \item \textbf{První část:} Lineární zobrazení jsou vyjádřena právě jednou unikátní maticí (bijekce mezi zobrazeními a maticemi), sčítáním zobrazení $f_A,f_B$
    získáváme zobrazení $f_{A+B}$, skalárním násobkem zobrazení $t\cdot f_A$ získáváme zobrazení $f_{tA}$.
    Z toho plyne, že prostor $Hom(\mathbb{R}^2,\mathbb{R}^2)$ je izomorfní s prostorem všech příslušných matic $\mathbb{R}^{2\times2}$.
    
    \textbf{Lemma} (*): Necht' libovolná báze B prostoru $\mathbb{R}^{2\times2}$. Prostor všech matic (A) $\mathbb{R}^{2\times2}$ je izomorfní
      s prostorem všech matic ($[A]^B_B$) $\mathbb{R}^{2\times2}$ vzhledem k bázi B.

    \begin{quote}
      \textbf{Důkaz} (*): Definujme zobrazení $f: \mathbb{R}^{2\times2} \to \mathbb{R}^{2\times2}$, které
        sobě přiřazuje příslušné matice $f(A) = [A]^B_B$.
        Ukažme, že je to izomorfismus (linearita+bijekce). Dle tvrzení 6.18 (vyjádření matice vzhledem k bázi) máme:
          \begin{align*}
            R &:= [id]^K_B \\
            [A]^B_B &= R^{-1} \cdot A \cdot R
          \end{align*}
        \begin{itemize}
          \item \textbf{Linearita}: Nechť $A_1,A_2 \in \mathbb{R}^{2\times2}$ a skalár $t\in\mathbb{R}$.
            Pak platí:
              \[
                f(A_1+A_2) = R^{-1} (A_1+A_2) R = R^{-1} (A_1R + A_2R) = R^{-1} A_1R + R^{-1} A_2R = f(A_1) + f(A_2)
              \]
              \[
                f(tA_1) = R^{-1} (tA_1) R = t(R^{-1} A_1 R) = t\cdot f(A_1)
              \]
            Tím je linearita splněna.
          \item \textbf{Bijekce}: Jelikož je matice přechodu $R$ regulární, lze rovnost výše takto upravit:
              \[
                [A]^B_B = R^{-1} \cdot A \cdot R \quad \Longrightarrow \quad R \cdot [A]^B_B  R^{-1} = A
              \]
            Z tohoto plyne, že každé matici $[A]^B_B$ lze najít vzor $A$. Tedy existuje inverz $f^{-1}(A)=R \cdot A \cdot R^{-1}$.
            $f$ je invertibilní, pak je bijektivní.
        \end{itemize}
        Jelikož je $f$ bijektivní a zachovává linearitu, jde o izomorfismus. $\quad \square$
    \end{quote}

    Protože platí lemma (*), pak z tranzitivity izomorfismu platí, že prostor $Hom(\mathbb{R}^2,\mathbb{R}^2)$ 
    je izomorfní s prostorem všech příslušných matic $\mathbb{R}^{2\times2}$ vzhledem k libovolné bázi B.

    Pak posloupnost $(f_1,f_2,f_3,f_4)$ je bází $Hom(\mathbb{R}^2,\mathbb{R}^2)$ právě tehdy, když
    je posloupnost $([f_1]^B_B,[f_2]^B_B,[f_3]^B_B,[f_4]^B_B)$ bází $\mathbb{R}^{2\times2}$.
    Ověřme, že posloupnost $([f_1]^B_B,[f_2]^B_B,[f_3]^B_B,[f_4]^B_B)$ je bází $\mathbb{R}^{2\times2}$:
    
    \begin{quote}
      Víme, že existuje izomorfismus $g: \mathbb{R}^{2\times2} \to \mathbb{R}^4$, $g\left(\left(\begin{array}{cc}
        a & b \\ c & d
      \end{array}\right)\right) = (a,b,c,d)^T$. Pak vyšetřeme, zda posloupnost
        \[
          \left(
            (1,1,0,1)^T, (1,1,1,0)^T, (1,0,1,1)^T, (0,1,1,1)^T
          \right)
        \]
      je bází $\mathbb{R}^4$. Posloupnost má 4 vektory, což je přesně dimenze $\mathbb{R}^4$,
      pak budou bází právě, když budou lineárně nezávislé.
      Položme vektory jako řádky do matice, matici zeliminujme, čímž vyšetříme nezávislost.

      \[
        \left(\begin{array}{cccc}
            1&1&0&1 \\
            1&1&1&0 \\
            1&0&1&1 \\
            0&1&1&1 \\
        \end{array}\right)
        \sim
        \left(\begin{array}{cccc}
            1&1&0&1 \\
            0&0&1&-1 \\
            0&-1&1&0 \\
            0&1&1&1 \\
        \end{array}\right)
         \sim
        \left(\begin{array}{cccc}
            1&1&0&1 \\
            0&1&1&1 \\
            0&0&1&-1 \\
            0&0&2&1 \\
        \end{array}\right)
         \sim
        \left(\begin{array}{cccc}
            1&1&0&1 \\
            0&1&1&1 \\
            0&0&1&-1 \\
            0&0&0&3 \\
        \end{array}\right)
      \]

      Vektory jsou nezávislé, takže jsou bází $\mathbb{R}^4$.

    \end{quote}

    Pak posloupnost zadaných matic je bází $\mathbb{R}^{2\times2}$, a tedy \underline{$(f_1,f_2,f_3,f_4)$ je bází $Hom(\mathbb{R}^2,\mathbb{R}^2)$}.

    \textbf{Druhá část:} Vyjádřeme matici $g$ vzhledem k bázi B.

    \begin{align*}
      R &:= [id]^B_C = \left(\begin{array}{cc}
        1.5 & 0.5 \\ -0.5 & 1.5
      \end{array}\right) = \left(\begin{array}{cc}
        0.6 & -0.2 \\ 0.2 & 0.6
      \end{array}\right)^{-1} \\
      [g]^B_B &\overset{6.18}{=} R^{-1} \cdot [g]^C_C \cdot R = \left(\begin{array}{cc}
        0.6 & -0.2 \\ 0.2 & 0.6
      \end{array}\right) \cdot \left(\begin{array}{cc}
        1&1 \\ 1&1       
      \end{array}\right) \cdot \left(\begin{array}{cc}
        1.5 & 0.5 \\ -0.5 & 1.5
      \end{array}\right) \\
      [g]^B_B &= \left(\begin{array}{cc}
        0.4 & 0.8 \\ 0.8 & 1.6
      \end{array}\right)
    \end{align*}

    Pro nalezení souřadnic opět využijme izomorfismus mezi $\mathbb{R}^{2\times2}$ a $\mathbb{R}^2$.

    \begin{align*}
      \left(\begin{array}{cccc|c}
          1&1&1&0 & 0.4\\
          1&1&0&1 & 0.8\\
          0&1&1&1 & 0.8\\
          1&0&1&1 & 1.6\\
      \end{array}\right)
      \sim
      \left(\begin{array}{cccc|c}
          1&0&0&0&0.4 \\
          0&1&0&0&-0.4 \\
          0&0&1&0&0.4 \\
          0&0&0&1&0.8 \\
      \end{array}\right)
    \end{align*}

    Souřadnice zobrazení $g$ v zadané bázi jsou aritmetický vektor \underline{\underline{$(0.4, -0.4, 0.4, 0.8)^T$}}.
\end{enumerate}

\end{document}
